% ---- pg-367 （书页 355） ----

\phantomsection\label{eq:8-35}
\begin{equation}
T=\frac{(n-1)}{(1-\bar{r})^2}\left[\sum_{i<k}\sum(r_{ik}-\bar{r})^2-\hat{\gamma}\sum_{k=1}^{p}(\bar{r}_k-\bar{r})^2\right]>\chi_{(p+1)(p-2)/2}^{2}(\alpha) \tag{8-35}
\end{equation}
则拒绝 $H_0$ 而接受 $H_1$, 其中 $\chi_{(p+1)(p-2)/2}^2(\alpha)$ 是自由度为 $(p+1)(p-2)/2$ 的 $\chi^2$ 分布的上 $100\alpha$ 百分位数.

\noindent\textbf{例 8.9}\quad (检验等相关结构)

从例 8.6, 由 $n=150$ 只雌性鼠生育后的体重构造的样本相关矩阵是
\[
\mathbf{R}=\begin{bmatrix}
1.0&0.750\ 1&0.632\ 9&0.636\ 3\\
0.750\ 1&1.0&0.692\ 5&0.738\ 6\\
0.632\ 9&0.692\ 5&1.0&0.662\ 5\\
0.636\ 3&0.738\ 6&0.662\ 5&1.0
\end{bmatrix}
\]
我们将用此相关矩阵说明在\hyperref[eq:8-35]{式(8-35)}中的大样本检验.

这里 $p=4$, 我们建立假设
\[
H_0:\boldsymbol{\rho}=\boldsymbol{\rho}_0=
\begin{bmatrix}
1&\rho&\rho&\rho\\
\rho&1&\rho&\rho\\
\rho&\rho&1&\rho\\
\rho&\rho&\rho&1
\end{bmatrix}
\]
\[
H_1:\boldsymbol{\rho}\neq\boldsymbol{\rho}_0
\]
利用\hyperref[eq:8-34]{式(8-34)}和\hyperref[eq:8-35]{式(8-35)}, 我们得到
\[
\bar{r}_1=\frac{1}{3}\times(0.750\ 1+0.632\ 9+0.636\ 3)=0.673\ 1,\qquad \bar{r}_2=0.727\ 1,
\]
\[
\bar{r}_3=0.662\ 6,\qquad \bar{r}_4=0.679\ 1
\]
\[
\begin{aligned}
\bar{r}&=\frac{2}{4\times 3}\times(0.750\ 1+0.632\ 9+0.636\ 3+0.692\ 5+0.738\ 6+0.662\ 5)\\
&=0.685\ 5
\end{aligned}
\]
\[
\begin{aligned}
\sum_{i<k}\sum(r_{ik}-\bar{r})^2&=(0.750\ 1-0.685\ 5)^2+(0.632\ 9-0.685\ 5)^2+\cdots+(0.662\ 5-0.685\ 5)^2\\
&=0.012\ 77
\end{aligned}
\]
\[
\sum_{k=1}^{p}(\bar{r}_k-\bar{r})^2=(0.673\ 1-0.685\ 5)^2+\cdots+(0.679\ 1-0.685\ 5)^2=0.002\ 45
\]
\[
\hat{\gamma}=\frac{(4-1)^2\times[1-(1-0.685\ 5)^2]}{4-(4-2)\times(1-0.685\ 5)^2}=2.132\ 9
\]
及
\[
T=\frac{(150-1)}{(1-0.685\ 5)^2}\times(0.012\ 77-2.132\ 9\times 0.002\ 45)=11.4
\]
既然 $(p+1)(p-2)/2=(5\times 2)/2=5$, 在\hyperref[eq:8-35]{式(8-35)}中, 检验的 5\% 临界值是 $\chi_5^2(0.05)=11.07$. 我们的检验统计量的值近似等于大样本 5\% 的临界点, 所以反对 $H_0$(等相关)的证据是有力的, 但还不能推翻 $H_0$.

同我们在例 8.6 见过的那样, 最小的特征值 $\hat{\lambda}_2,\hat{\lambda}_3$ 及 $\hat{\lambda}_4$ 相差细微, 而 $\hat{\lambda}_4$ 是比其他两个更小一点. 因此, 由于本问题里的大样本容量, 来自等相关结构的细小差别在统计上是显著的. \hfill$\blacksquare$
